% 1_introduction.tex
\section{Introduction}

Multi-agent systems built from large language models (LLMs) are
increasingly evaluated by having another LLM read the system's output
and score it. In multi-agent \emph{discussion} systems---several
role-conditioned agents deliberating on an open question---the scored
artifact is the discussion transcript itself, and the scores routinely
carry scientific weight: a recent open-source roundtable framework
reports that its decentralized speaking mechanism outperforms
single-model role-play (6.20 vs.\ 3.53 on a five-dimension judge scale)
and ranks five system components by their measured contribution
(host summary $-1.20$, confidence arbitration $-0.92$, jitter $-0.80$,
iron rules $-0.53$, opening $-0.40$)~\cite{whospeaks2024}. Conclusions
of this form---``this component matters more than that one''---now
appear across the multi-agent literature.

All such conclusions inherit an unexamined dependency: the identity of
the judge model. LLM judges are known to be imperfect measuring
instruments in general~\cite{zheng2023judging,wang2024canyoutrust},
and large-scale audits have shown that judge rankings move across
benchmarks and that reliability is item-specific rather than
model-level~\cite{reliability2026,allverdicts2026}. What has not been
measured is the consequence for the specific inferential move that
multi-agent papers make: a single judge scores transcripts of system
variants, and the resulting differences are read as properties of the
\emph{system}. If a different judge reverses the ordering of the
variants, the component ranking is not established as a property of the
system.

This paper is an audit of that inferential move, conducted on a system
whose source, protocol, and original evaluation are fully available to
us. We ask three research questions. \textbf{RQ1:} How consistent are
different judges' quality assessments of the same multi-agent
discussions? \textbf{RQ2:} Through what mechanism does judge identity
act---a preference for outputs of the judge's own family, or a general
leniency offset? \textbf{RQ3:} Do the component-importance conclusions
of a single-judge evaluation replicate when the judge changes?

Our approach has three parts. First, we re-implemented the audited
system's speaking protocol as a headless testbed with no third-party
dependencies (system prompt, decision schema, confidence arbitration,
jitter algorithm, and host prompts reproduced from the source;
deviations listed in \S\ref{sec:method}), so that every condition
differs only in the factor under study. Second, we ran a $2\times2$
factorial experiment crossing driver model with judge model, plus a
replication of the system's component ablation under multiple judges,
including a same-machine replication in which only the judge varies.
Third, we scored all transcripts with three judges from two
vendors---two models of the driver's vendor and one cross-vendor
model---so that judge identity is a controlled factor rather than an
assumption.

\paragraph{Contributions.}
(i)~A protocol-faithful, released testbed for multi-agent discussion
evaluation (all prompts and per-round decision logs included).
(ii)~A quantification of cross-judge agreement in this setting: rank
correlation is low and heterogeneous ($\rho\leq+0.28$ among the three
analyzed judges, $+0.47$ at most including a fourth released judge; one
cross-vendor pair reaches $+0.61$ in a parser-clean sensitivity
analysis), and level differences reach 0.87 points on a 10-point
scale.
(iii)~A mechanism decomposition showing that the dominant inter-judge
effect is a leniency offset whose confidence interval excludes zero and
that persists on transcripts the judge did not produce, while
self-preference is not detected---refining the standard mitigation
advice that a different judge family suffices.
(iv)~A replication audit showing that a single judge's
component-importance conclusions vary with the judge (Cliff's $\delta$
from $-1.00$ to $+0.78$ across judges and drivers), that no judge
reproduces the audited evaluation's component ordering in the
same-machine replication, and that no judge reproduces the audited
effect magnitudes; with three-to-five repetitions per condition, no
contrast survives multiplicity correction, so we report direction and
effect size rather than significance.

All quantitative claims are scoped to the studied setting: one
discussion framework, one ablation topic, five repetitions per ablation
condition (three in the same-machine replication), and 32 discussions
in total; we state these boundaries where they apply.

\paragraph{Scope and posture.} The first author is the developer of the
audited framework, and the audited evaluation is the first author's own
prior work. We treat this as a case study of an evaluation protocol,
not as evidence about the framework's quality: mechanism metrics
(\S\ref{sec:mechanism}) indicate the framework operates as designed, and
nothing here claims otherwise.

The rest of the paper is organized as follows. \S\ref{sec:related}
positions the work against the LLM-judge reliability literature,
multi-agent evaluation, and reproducibility audits. \S\ref{sec:method}
describes the audited protocol, the testbed, and the experimental
design. \S\ref{sec:results} reports the three findings.
\S\ref{sec:discussion} states the minimum reporting norm we propose,
and \S\ref{sec:conclusion} concludes.
